\section{Gated Recurrent Unit (GRU)}

The Python part of this exercise is available in \myhref{https://github.com/LosDelDGIIM/LosDelDGIIM.github.io/blob/main/subjects/Erasmus-DUE/GKI/Exercises/Exercise9.ipynb}{this Jupyter Notebook}.\\


\begin{ejercicio}
    Describe the main components of a GRU cell and their function.

    A GRU cell consists of three main components:
    \begin{itemize}
        \item \ul{Update Gate}: This gate determines how much of the previous hidden state should be retained and how much of the new candidate hidden state should be used. It helps the model to decide whether to keep the old information or to update it with new information.
        \item \ul{Reset Gate}: This gate controls how much of the previous hidden state should be ignored when calculating the new candidate hidden state. It allows the model to reset its memory when necessary, which can help in capturing short-term dependencies.
        \item \ul{Candidate Hidden State}: This is a potential new hidden state that is computed based on the current input and the previous hidden state, modulated by the reset gate.
    \end{itemize}
    The final hidden state is then computed as a combination of the previous hidden state and the candidate hidden state, weighted by the update gate.
\end{ejercicio}
\begin{ejercicio}
    How do the update and reset gates in a GRU cell differ in their operation?

    The update gate in a GRU cell controls the extent to which the previous hidden state is retained versus how much of the new candidate hidden state is used. It essentially decides whether to keep the old information or to update it with new information. In contrast, the reset gate determines how much of the previous hidden state should be ignored when calculating the new candidate hidden state. It allows the model to reset its memory when necessary, which can help in capturing short-term dependencies. In summary, the update gate manages the flow of information from the past to the future, while the reset gate manages how much of the past information is considered when generating new candidate states.
\end{ejercicio}
\begin{ejercicio}
    Name two advantages of GRUs over simple RNNs.

    Two advantages of GRUs over simple RNNs (which are really related) are:
    \begin{itemize}
        \item \ul{Mitigation of the Vanishing Gradient Problem}: GRUs, like LSTMs, use gating mechanisms that help to preserve gradients during backpropagation.
        \item \ul{Learning Long-Term Dependencies}: GRUs can capture long-term dependencies in sequential data more effectively than simple RNNs, which struggle with this due to the vanishing gradient problem.
    \end{itemize}
\end{ejercicio}
\begin{ejercicio}
    Compare GRUs and LSTMs with regard to their structure and complexity.

    GRUs and LSTMs are both types of recurrent neural network architectures designed to address the vanishing gradient problem. However, they differ in their structure and complexity:
    \begin{itemize}
        \item \ul{Structure}: LSTMs have three gates (input, forget, and output gates) and a cell state, while GRUs have only two gates (update and reset gates). This makes GRUs simpler in terms of architecture compared to LSTMs.
        \item \ul{Complexity}: Due to their simpler structure, GRUs typically have fewer parameters than LSTMs, which can lead to faster training and inference times. However, the choice between GRUs and LSTMs often depends on the specific task and dataset, as LSTMs may perform better on certain tasks that require modeling more complex dependencies.
    \end{itemize}
\end{ejercicio}
\begin{ejercicio}
    Give a typical application example for GRUs.

    A typical application example for GRUs is in natural language processing (NLP) tasks, such as:
    \begin{itemize}
        \item \ul{Machine Translation}: GRUs can be used in sequence-to-sequence models for translating text from one language to another, where they help to capture the dependencies between words in a sentence.
        \item \ul{Speech Recognition}: GRUs can be employed in models that convert spoken language into text, as they can effectively model the temporal dependencies in audio data.
        \item \ul{Text Generation}: GRUs can be used in language models to generate coherent and contextually relevant text based on a given input prompt.
    \end{itemize}
\end{ejercicio}
\begin{ejercicio}
    How do GRUs help to reduce the vanishing gradient problem?

    GRUs help to reduce the vanishing gradient problem through their gating mechanisms, specifically the update and reset gates. The update gate allows the model to retain information from previous time steps, which helps to preserve gradients during backpropagation. The reset gate allows the model to ignore irrelevant past information when generating new candidate hidden states, which can also help to maintain a stronger gradient signal. By controlling the flow of information and allowing for long-term dependencies, GRUs mitigate the vanishing gradient problem that simple RNNs often face when trying to learn from long sequences.
\end{ejercicio}